\subsection{ORBIT TYPES}

Let \(G\) be a topological group operating continuously on a separated topological space \(E\) . For every point \(x\) of \(E\) , denote by \(G_x\) the fixer of \(x\) in \(G\) , and assume that the canonical map \(G / G_x \to Gx\) is a homeomorphism; this is notably the case in the following two situations:

\begin{enumerate}
\def\labelenumi{\alph{enumi})}
\item
  the topologies of G and E are discrete;
\item
  G operates properly on E (General Topology, Chap. III, §4, no. 2, Prop. 4), for example, G is compact (General Topology, Chap. III, §4, no. 1, Prop. 2).
\end{enumerate}

Denote by \(\mathcal{T}\) the set of conjugacy classes of closed subgroups of G. For every \(x\in \mathrm{E}\) , we call the orbit type of \(x\) , or sometimes the type of \(x\) , the class of \(\mathrm{G}_x\) in \(\mathcal{T}\) ; two points of the same orbit are of the same orbit type (Algebra, Chap. I, §5, no. 2, Prop. 2); two orbits are of the same type if and only if they are isomorphic as G-sets (Algebra, Chap. I, §5, no. 5, Th. 1). For every \(t\in \mathcal{T}\) , denote by \(\mathrm{E}_{(t)}\) the set of points of E of type \(t\) , that is, the union of the orbits of type \(t\) ; this is a stable subset of E. For \(\mathrm{H}\in t\) , we also write \(\mathrm{E}_{(\mathrm{H})}\) for \(\mathrm{E}_{(t)}\) ; for example, \(\mathrm{E}_{(\mathrm{G})}\) is the closed subspace of E consisting of the points fixed by G.

Give T the following preorder relation:

\(t \leq t'\Longleftrightarrow\) there exists \(\mathrm{H} \in t\) and \(\mathrm{H}' \in t'\) such that \(\mathrm{H} \supset \mathrm{H}'\) .

Let \(\Omega\) and \(\Omega'\) be two orbits of G on E, \(t\) and \(t'\) their types; then \(t \leq t'\) if and only if there exists a G-morphism (necessarily surjective and continuous) from \(\Omega'\) to \(\Omega\) .

Let \(x, x'\) be in E, and \(t, t'\) their types; then \(t \leq t'\) if and only if there exists \(a \in G\) such that \(aG_{x'}a^{-1} \subset G_x\) .

Lemma 6. Let G be a Lie group.

\begin{enumerate}
\def\labelenumi{\alph{enumi})}
\item
  Every decreasing sequence of compact subgroups of \(\mathbf{G}\) is stationary.
\item
  Let H and \(H'\) be two compact subgroups of G such that \(H \subset H'\) and such that there exists an isomorphism (of topological groups) from \(H'\) to H. Then \(H = H'\) .
\item
  With the relation \(t \leq t'\) , the set \(\mathcal{T}\) is a noetherian ordered set (Theory of Sets, Chap. III, §6, no. 5, text preceding Prop. 7).
\item
  Let \((\mathrm{H}_{i})_{i\geq1}\) be a decreasing sequence of compact subgroups of G; these are Lie subgroups of G (Chap. III, §8, no. 2, Th. 2). The sequence of integers \((\dim\mathrm{H}_{i})_{i\geq1}\) is decreasing, hence stationary, so there exists an integer N such that the subgroups \(H_{i}\) have the same identity component for \(i\geq N\) . Then the decreasing sequence of positive integers \((\mathrm{H}_{i}:(\mathrm{H}_{i})_{0})_{i\geq N}\) is stationary, so \(H_{i}=H_{i+1}\) for i sufficiently large.
\item
  Let \(f\) be an isomorphism from \(\mathrm{H}'\) to \(\mathrm{H}\) . The sequence \((f^n(\mathrm{H}))_{n\geq 0}\) is a decreasing sequence of compact subgroups of \(\mathrm{G}\) , so \(f^n(\mathrm{H}) = f^{n + 1}(\mathrm{H})\) for \(n\) sufficiently large, by \(a\) ). Since \(f\) is an isomorphism, this implies that \(f(\mathrm{H}) = \mathrm{H} = f(\mathrm{H}')\) , so \(\mathrm{H} = \mathrm{H}'\) .
\item
  Let \(t, t' \in \mathcal{T}\) be such that \(t \leq t'\) and \(t' \leq t\) . Then, there exist \(\mathrm{H}, \mathrm{H}_1 \in t\) and \(\mathrm{H}', \mathrm{H}_1' \in t'\) such that \(\mathrm{H} \supset \mathrm{H}'\) and \(\mathrm{H}_1 \subset \mathrm{H}_1'\) . Let \(g\) and \(g'\) be two elements of \(G\) such that \(\mathrm{H}_1 = g\mathrm{H}g^{-1}\) and \(\mathrm{H}_1' = g'\mathrm{H}'g'^{-1}\) ; put \(u = g'^{-1}g\) . Then
\end{enumerate}

\[
u \mathrm{H} u ^ {- 1} \subset \mathrm{H} ^ {\prime} \subset \mathrm{H};
\]

by \(b\) ), this implies that \(u\mathrm{H}u^{-1} = \mathrm{H}\) , so \(\mathrm{H}' = \mathrm{H}\) and \(t' = t\) . Thus, the set \(\mathcal{T}\) is ordered, and it is noetherian by \(a\) ).

THEOREM 2. Let \(\mathbf{G}\) be a Lie group operating properly on \(\mathbf{X}\) , such that the law of operation \((g,x)\mapsto gx\) is of class \(\mathbf{C}^r\) . Assume that \(\mathbf{X}\) is paracompact.

\begin{enumerate}
\def\labelenumi{\alph{enumi})}
\item
  The map which associates to any point of X its orbit type has the following semi-continuity property: let \(x \in X\) and let \(t \in T\) be its orbit type; there exists a stable open neighbourhood U of x such that, for any \(u \in U\) , the type of u is \(\geq t\) .
\item
  For all \(t \in \mathcal{T}\) , \(\mathrm{X}_{(t)}\) is a submanifold of \(\mathrm{X}\) , the equivalence relation on \(\mathrm{X}_{(t)}\) induced by the operation of \(\mathrm{G}\) is regular (Differentiable and Analytic Manifolds, Results, 5.9.5), and the morphism \(\mathrm{X}_{(t)} \to \mathrm{X}_{(t)} / \mathrm{G}\) is a bundle.
\item
  Assume that X/G is connected. Then the set of orbit types of elements of X has a largest element \(\tau\) ; moreover, \(\mathrm{X}_{(\tau)}\) is a dense open subset of X and \(\mathrm{X}_{(\tau)}/\mathrm{G}\) is connected.
\end{enumerate}

Let \(x\) be a point of \(\mathrm{X}\) and \(t \in \mathcal{T}\) its type. To prove \(a)\) and \(b)\) , we can replace \(\mathrm{X}\) by a stable open set containing \(x\) , and hence (Prop. 6) can assume that \(\mathrm{X}\) is of the form \(\mathrm{G} \times^{\mathrm{H}}\mathrm{W}\) , where \(\mathrm{W}\) is the space of a finite dimensional analytic linear representation of a compact subgroup \(\mathrm{H}\) of \(\mathrm{G}\) , the point \(x\) being the image \(p(e,0)\) of \((e,0) \in \mathrm{G} \times \mathrm{W}\) under the canonical projection \(p: \mathrm{G} \times \mathrm{W} \to \mathrm{G} \times^{\mathrm{H}}\mathrm{W}\) . If \(u = p(g,y) \in \mathrm{G} \times^{\mathrm{H}}\mathrm{W}\) and \(a \in \mathrm{G}\) , then \(au = u\) if and only if there exists \(h \in \mathrm{H}\) with \((ag,y) = (gh^{-1},hy)\) , that is, if \(a \in g\mathrm{H}_yg^{-1}\) . Thus, \(\mathrm{G}_u = g\mathrm{H}_yg^{-1}\) ; in particular, \(\mathrm{G}_x = \mathrm{H}\) , so \(\mathrm{G}_u\) is conjugate to a subgroup of \(\mathrm{G}_x\) , which proves that the type of \(u\) is \(\geq t\) , hence \(a)\) .

Moreover, u is of type t if and only if \(G_{u}\) is conjugate to H in G, or equivalently that \(H_{y}\) is conjugate to H in G; by Lemma 6 b), this means that \(H_{y} = H\) , and hence that y is fixed by H. If \(W'\) is the vector subspace of W consisting of the elements fixed by H, it follows that \(\mathrm{X}_{(t)}\) can be identified with \(G \times {}^{H}W'\) , and hence also with \(G/H \times W'\) , hence b).

To prove c), observe that the assumption that X/G is connected implies that X is pure of finite dimension: indeed, for all \(k \geq 0\) , denote by \(X_{k}\) the set of points \(x \in X\) such that \(\dim_{x} X = k\) ; then \(X_{k}\) is open and closed in X, and stable under G, so X is equal to one of the \(X_{k}\) .

We now prove c) by induction on the dimension of X, the assertion being clear for \(\dim X = 0\) . Let \(\tau\) be a maximal element among the orbit types of the points of X (such an element exists by Lemma 6 c)). We shall prove the following:

\(c^{\prime})\) For every subset A of \(\mathrm{X}_{(t)}\) , open and closed in \(\mathrm{X}_{(\tau)}\) and stable under G, the closure \(\overline{\mathrm{A}}\) of A in X is open.

This assertion implies \(c\) ). Indeed, note first that \(\mathrm{X}_{(\tau)}\) is open in \(\mathrm{X}\) , by \(a\) ; assertion \(c'\) ) implies that \(\overline{\mathrm{X}}_{(\tau)}\) is open and closed in \(\mathrm{X}\) , hence equal to \(\mathrm{X}\) since it is stable under \(\mathrm{G}\) and \(\mathrm{X} / \mathrm{G}\) is connected. Let \(\mathrm{A}\) be a non-empty open and closed subset of \(\mathrm{X}_{(\tau)}\) stable under \(\mathrm{G}\) ; by \(c'\) ), \(\overline{\mathrm{A}}\) is open and closed in \(\mathrm{X}\) and stable under \(\mathrm{G}\) , hence equal to \(\mathrm{X}\) ; this implies that \(\mathrm{A}\) is dense in

\(X_{(\tau)}\) , hence equal to \(X_{(\tau)}\) . Consequently, every non-empty open and closed subset of \(X_{(\tau)}/G\) is equal to \(X_{(\tau)}/G\) , which proves that \(X_{(\tau)}/G\) is connected. Finally, since \(X_{(\tau)}\) is dense in X, it follows from a) that every point of X is of type \(\leq \tau\) ; in other words, \(\tau\) is the largest element among the orbit types of the points of X.

We now prove \(c'\) . We can assume that A is non-empty; let \(x \in \overline{A}\) . It suffices to prove that \(\overline{A}\) is a neighbourhood of x. For this we can, as above, assume that \(X = G \times {}^{H}W\) with H compact, x being the canonical image of \((e,0)\) . Assume first that H operates trivially on W: then X can be identified with \((G/H) \times W\) , and \(X_{(\tau)}/G = X/G\) is homeomorphic to W, hence connected; thus, A/G = X/G, so A = X. Assume from now on that H does not operate trivially on W. Choose a scalar product on W invariant under the compact group H; let S be the unit sphere in W (the set of vectors of norm 1). Note that S/H is connected: indeed, if \(\dim(W) \geq 2\) , S is connected, and if \(\dim(W) = 1\) , S is a space of two points on which H operates non-trivially. Put \(Y = G \times {}^{H}S\) ; this is a closed submanifold of X, stable under G, of codimension 1, and Y/G, which is homeomorphic to S/H, is connected. Thus, by the induction hypothesis, there exists a maximal orbit type \(\theta\) for Y, the set \(Y_{(\theta)}\) is open and dense in Y, and \(Y_{(\theta)}/G\) is connected.

Consider the operation of \(R_{+}^{*}\) on X induced by passage to the quotient by the law of operation \((\lambda, (g, w)) \mapsto (g, \lambda w)\) of \(R_{+}^{*}\) on G × W. Two points of X conjugate under this operation are of the same orbit type; consequently, \(\mathrm{X}_{(\theta)}\) contains \(\mathbf{R}_{+}^{*}\mathrm{Y}_{(\theta)}\) , which is a dense open subset of X. But \(\mathrm{X}_{(\tau)}\) is open by a), and hence meets \(\mathrm{X}_{(\theta)}\) , so \(\theta = \tau\) .

On the other hand, the homeomorphism \((\lambda,w)\mapsto\lambda w\) from \(R_{+}^{*}\times S\) to \(W-\{0\}\) (General Topology, Chap. VI, §2, no. 3, Prop. 3) induces a homeomorphism from \(R_{+}^{*}\times(S/H)\) to \((\mathbf{R}_{+}^{*}\mathrm{S})/\mathrm{H}\) , hence also from \(R_{+}^{*}\times(Y/G)\) to \((\mathbf{R}_{+}^{*}\mathrm{Y})/\mathrm{G}\) , and from \(R_{+}^{*}\times(\mathrm{Y}_{(\theta)}/\mathrm{G})\) to \((\mathbf{R}_{+}^{*}\mathrm{Y}_{(\theta)})/\mathrm{G}\) . Thus, \((\mathbf{R}_{+}^{*}\mathrm{Y}_{(\theta)})/\mathrm{G}\) is connected, and \(X_{(\tau)}/\mathrm{G}\) , which contains a connected dense subset, is itself connected (General Topology, Chap. I, §11, no. 1, Prop. 1). Consequently, A is equal to \(X_{(\tau)}\) , hence is dense in X, and \(\overline{A}\) is a neighbourhood of x. This completes the proof of the theorem.

With the notations in Th. 2 c), the points of \(\mathrm{X}_{(\tau)}\) are said to be principal and their orbits are called principal orbits. If x is a point of X, and if \(G \times G_{x}W\) is a linear tube in X around the orbit of x, the point x is principal if and only if \(G_{x}\) operates trivially on W, that is, if the tube is of the form \((\mathrm{G}/\mathrm{G}_{x}) \times \mathrm{W}\) .

Examples. 1) Let G be a connected compact Lie group, operating on itself by inner automorphisms. The fixer of an element x of G is simply the centralizer \(Z(x)\) of x in G; it contains every maximal torus containing x. It follows that the largest orbit type \(\tau\) is the conjugacy class of the maximal tori of G. The open set \(\mathrm{G}_{(\tau)}\) is the set of very regular elements of G ( \(\S5\) , no. 1, Remark). Assume that G is simply-connected. Then \(\mathrm{G}_{(\tau)}\) is equal to the set \(G_{r}\) of regular elements of G ( \(\S5\) , no. 2, Remark 2); if A is an alcove of a Cartan subalgebra t of \(\mathfrak{g} = \mathrm{L}(\mathrm{G})\) , the composite map \(\pi: A \xrightarrow{\exp} G_{r} \longrightarrow G_{r}/\mathrm{Int}(\mathrm{G})\) is an isomorphism of analytic manifolds. Indeed, this is a homeomorphism (§5, no. 2, Cor. of Prop. 2); let \(a \in A\) , put \(t = \exp a\) and identify \(T_{t}(\mathrm{G})\) with g by means of the translation \(\gamma(t)\) . The tangent map \(T_{a}(\pi)\) can then be identified with the composite of the canonical injection \(t \to g\) and the quotient map \(g \to g/\mathrm{Im}(\mathrm{Ad} t^{-1} - 1)\) . Since t is regular, \(T_{a}(\pi)\) is an isomorphism, hence the stated result (Differentiable and Analytic Manifolds, Results, 5.7.8).

\begin{enumerate}
\def\labelenumi{\arabic{enumi})}
\setcounter{enumi}{1}
\tightlist
\item
  Let E be a real euclidean affine space, H a set of hyperplanes of E, W the group of displacements of E generated by the orthogonal reflections with respect to the hyperplanes of H. Assume that H is stable under W and that the group W, with the discrete topology, operates properly on E.
\end{enumerate}

The preceding can be applied to the operation of W on E. The fixer of a point x of E is the subgroup of W generated by the reflections with respect to the hyperplanes of H containing x (Chap. V, §3, no. 3, Prop. 2). Consequently, the largest orbit type \(\tau\) is the class of the subgroup \(\{Id_{E}\}\) , and \(\mathrm{E}_{(\tau)}\) is the union of the chambers of E. Note that in this case the covering \(\mathrm{E}_{(\tau)} \to \mathrm{E}_{(\tau)}/\mathrm{W}\) is trivial, and in particular \(\mathrm{E}_{(\tau)}\) is not connected if H is non-empty.

